\chapter{Abbreviations}
\newlength{\taiubh}
\settowidth{\taiubh}{TaiUBh} % looks like this is the widest abbreviation
\addtolength{\taiubh}{12pt}
\newlength{\rightcolumn}
\setlength{\rightcolumn}{\textwidth}
\addtolength{\rightcolumn}{-\taiubh}

\newcommand{\abbrindex}[4]{%
  \index{#1@#2|see{#3}}%
  #2 & #3 \citep{#4} \\[0.50ex]
}
% \index{BAU@BĀU|see{Bṛhadāraṇyakopaniṣad}}
% \index{BAUBh@BĀUBh|see{Bṛhadāraṇyako\-pa\-ni\-ṣa\-d\-bhā\-ṣya}}
\section{Texts and editions}
\setlength\LTleft{0pt}
\setlength\LTright{0pt}
\begin{longtable}{@{}l p{\rightcolumn}@{}}
\abbrindex{AA}{AA}{Aṣṭādhyāyī of Pāṇini}{vasu1962ashtadhyayi} %vasu...
% ANSS
\abbrindex{BAU}{BĀU}{Bṛhadāraṇyako\-pa\-ni\-ṣa\-d}{bau}
\abbrindex{BAUBh}{BĀUBh}{Bṛ\-ha\-dā\-ra\-ṇya\-ko\-pa\-ni\-ṣa\-d\-bhā\-ṣya ascribed to Śaṅkara}{bau}
\abbrindex{BhG}{BhG}{Bhagavadgītā}{bhgbh}
\abbrindex{BhGBh}{BhGBh}{Bhagavadgītābhāṣya ascribed to Śaṅkara}{bhgbh}
\abbrindex{BS}{BS}{Brahmasūtra}{bsbh}
\abbrindex{BSBh}{BSBh}{Bra\-hma\-sū\-tra\-śā\-ṅka\-ra\-bhā\-ṣya of Śaṅkara}{bsbh}
\abbrindex{BSi}{BSi}{Brahmasiddhi of Maṇḍana Miśra}{bs}
\abbrindex{ChU}{ChU}{Chāndogyopaniṣad}{upa2}
\abbrindex{ChUBh}{ChUBh}{Chā\-ndo\-gyo\-pa\-ni\-ṣa\-dbhā\-ṣya ascribed to Śaṅka\-ra}{upa2}
% CSS
% \abbrindex{IUBh}{ĪUBh}{Īśopaniṣadbhāṣya ascribed to Śaṅkara}{} 
% \abbrindex{IU}{ĪU}{Īśopaniṣad}{}
\abbrindex{JS}{JS}{Jaiminisūtra a.k.a.\ Mī\-māṃ\-sā\-sū\-tra}{frauwallner:karma,js1,js1a,js2,js3,js4,js5,js6,js7}
\abbrindex{KAS}{KAŚ}{Kauṭilyārthaśāstra}{as:kangle_ed1}
\abbrindex{KaU}{KaU}{Kaṭhopaniṣad}{upa1} 
\abbrindex{KaUBh}{KaUBh}{Kaṭhopaniṣadbhāṣya ascribed to Śaṅkara}{upa1} 
\abbrindex{KeU}{KeU}{Kenopaniṣad}{upa1} 
\abbrindex{KeUBh}{KeUBh}{Kenopaniṣadbhāṣya ascribed to Śaṅkara}{upa1} 
% KU
% KUPBh
\abbrindex{Manu}{Manu}{Manusmṛti}{olivelle:manu_ed}% Olivelle 2006 
\abbrindex{MaU}{MāU}{Māṇḍūkyopaniṣad}{upa1}
\abbrindex{MaUBh}{MāUBh}{Mā\-ṇḍū\-kyo\-pa\-ni\-ṣa\-dbhā\-ṣya ascribed to Śa\-ṅka\-ra}{upa1}
\abbrindex{MBh}{MBh}{Mahābhārata}{poona:mbh} 
\abbrindex{MMK}{MMK}{Mū\-la\-ma\-dhya\-mi\-ka\-kā\-ri\-kās of Nāgārjuna}{mmk} 
% \abbrindex{MS}{MS}{Manusmṛti???}{} 
\abbrindex{MuU}{MuU}{Muṇḍakopaniṣad}{upa1}
\abbrindex{MuUBh}{MuUBh}{Mu\-ṇḍa\-ko\-pa\-ni\-ṣa\-dbhā\-ṣya ascribed to Śaṅkara}{upa1} 
\abbrindex{NBh}{NBh}{Nyāyabhāṣya}{ns}
\abbrindex{NBhu}{NBhū}{Nyāyabhūṣaṇa of Bhā\-sa\-rva\-jña}{nbhusana} 
% \abbrindex{NK}{NK}{Nyāya\ldots???}{}
\abbrindex{NKd}{NKd}{Nyāyakandalī}{pdhs}
\abbrindex{Nkn}{NKṇ}{Nyāyakaṇikā of Vācaspati Miśra}{vidhiviveka,vidhiviveka:tara,vidhiv:stern}
\abbrindex{NM}{NM}{Nyāyamañjarī of Bhaṭṭa Jayanta}{nm:kss}
\abbrindex{NS}{NS}{Nyāyasūtra}{ns}
\abbrindex{NV}{NV}{Nyāyavārttika}{ns}
\abbrindex{NVTT}{NVTṬ}{Nyā\-ya\-vā\-rtti\-ka\-tā\-tpa\-rya\-ṭī\-kā of Vācaspati Miśra}{ns}
\abbrindex{PDhS}{PDhS}{Padārthadharmasaṅgraha}{pdhs}
% \abbrindex{PP}{PP}{Prakaraṇapañcikā of Śālikanātha}{pp}
\abbrindex{PaSu}{PāSū}{Pāśupatasūtra}{pasupatasutra}
\abbrindex{PraPa}{PraPa}{Prakaraṇapañcikā}{pp}
\abbrindex{PYS}{PYŚ}{Pātañjalayogaśāstra}{ys:anss,yvi}
\abbrindex{PV}{PV}{Pramāṇavārttika of Dha\-rma\-kīrti}{pvv,pv:pandeya}
\abbrindex{SBh}{ŚBh}{Śābarabhāṣya}{frauwallner:karma,js1,js1a,js2,js3,js4,js5,js6,js7}
\abbrindex{SK}{SK}{Sāṃkhyakārikās}{yd} %yd
\abbrindex{SRTS}{ŚRTS}{Tattvasaṃgraha of Śā\-nta\-ra\-kṣita}{ts:bbh}
\abbrindex{SS}{SS}{Sphoṭasiddhi of Maṇḍana Miśra}{SS}
\abbrindex{SV}{ŚV}{Ślokavārttika of Kumārila}{sv:umbeka,sv:partha,svkasika} 
\abbrindex{TaiU}{TaiU}{Taittirīyopaniṣad}{upa1}
\abbrindex{TaiUBh}{TaiUBh}{Tai\-tti\-rī\-yo\-pa\-ni\-ṣa\-dbhā\-ṣya ascribed to Śaṅkara}{upa1}
\abbrindex{TS}{TS}{Taittirīyasaṃhitā}{weber:tais}
\abbrindex{TSP}{TSP}{Tattvasaṃgrahapañjikā of Kamalaśīla}{ts:bbh}
%\abbrindex{TP}{TP}{???}{}
\abbrindex{TV}{TV}{Tattvavaiśāradī (a commentary on the Yo\-ga\-bhā\-ṣya) of Vā\-ca\-s\-pa\-ti Miśra}{ys:anss,yvi}
\abbrindex{Upad}{Upad}{Upadeśasāhasrī}{mayeda:upad}
\abbrindex{VidhiV}{VidhiV}{Vidhiviveka of Maṇḍana Miśra}{vidhiviveka,vidhiviveka:tara,vidhiv:stern}
\abbrindex{VMBh}{VMBh}{Mahābhāṣya of Patañjali}{mbh1,mbh3}
\abbrindex{VP}{VP}{Vākyapadīya of Bhartṛhari}{rau:text,rau:index,iyer:vp2text} 
\abbrindex{VS}{VS}{Vaiśeṣikasūtra}{vs}
\abbrindex{YBh}{YBh}{Yogabhāṣya}{ys:anss,yvi,maas:2006}
\abbrindex{YD}{YD}{Yuktidīpikā}{wezler:yd}
\abbrindex{YS}{YS}{Yogasūtra}{ys:anss,yvi,maas:2006}
\abbrindex{YV}{YV}{Yogavārttika of Vi\-jñā\-na\-bhi\-kṣu}{rukmani:yvarttika}
\abbrindex{YVi}{YVi}{Pātañjalayogaśāstravivaraṇa}{yvi}
\index{Jayamangala@Jayamaṅgalā|seealso{Sāṃkhyajayamaṅgalā}}
\index{Mimamsa@Mīmāṃsā|seealso{Mīmāṃsaka}}
\index{Samkhyajayamangala@Sāṃkhyajayamaṅgalā|seealso{Jayamaṅgalā}}
% \abbrindex{}{}
% \abbrindex{}{}
% \abbrindex{}{}
% \abbrindex{}{}
% \abbrindex{}{}
% \abbrindex{}{}
% \abbrindex{}{}
% \abbrindex{}{}
% \abbrindex{}{}
% \abbrindex{}{}
% \abbrindex{}{}
\end{longtable}

\section{Sigla and Others}
See pp.~\pageref{sec:symbols} for details.
\setlength\LTleft{0pt}
\setlength\LTright{0pt}
\begin{longtable}{@{}ll@{}}
	MS & manuscript \\
	em. & emendation \\
	conj. & conjecture \\
	\Tm & Trivandrum manuscript \\
	\Tmac & ditto before correction \\
	\Tmpc & ditto after correction \\
	\Td & A transcript of the Trivandrum manuscript \\
	\Tdac & ditto before correction \\
	\Tdpc & ditto after correction \\
	T & A reading common in \Tm\ and \Td \\
	L & Lahore manuscript \\
	\Lac & ditto before correction \\
	\Lpc & ditto after correction \\
	M & Madras Government Oriental Manuscript Library manuscript \\
	\Mac & ditto before correction \\
	\Mpc & ditto after correction \\
	A & Adyar Library manuscript \\
	\Aac & ditto before correction \\
	\Apc & ditto after correction \\
	\Me & 1952 edition of the YVi \\
	$\Sigma$ & all the witnesses, including the 1952 edition \\
	\mms & the following text is lost in \Tm \\
	\kern 5.3pt \mme & the preceding text is lost in \Tm\\
	\raisebox{-0.5ex}{\mds} & the following text is reported missing in \Td \\
	\raisebox{-0.5ex}{\kern 5.3pt \mde} & the preceding text is reported missing in \Td \\
	\dag & the enclosed text is suspected to be corrupt \\
	\begin{uncertain}{\dn क}\end{uncertain} & (text in a lighter shade) uncertain text \\
	\ldots & in critical text, suspected loss of text \\[0pt]
	[\mist{1}] & physically lost sign in a manuscript \\
	\mist{1} & a sign indicated as lost in the exemplar \\
	\hlg{{\dn क}} & highly illegible but ascertainable as {\dn क} \\
	\ilg & an illegible (but not lost) sign \\
	\ccs {\dn क}\cce & cancelled {\dn क} \\
	\ins{\dn क}\ine & inserted {\dn क} \\
	\open & space left by scribe
\end{longtable}

% \begin{description}
% 	\item[MS]{manuscript}
% 	\item[em.]{emendation}
% 	\item[conj.]{conjecture}
% 	\item[\Tm]{Trivandrum manuscript}
% 	\item[\Tmac]{Trivandrum manuscript before correction}
% 	\item[\Tmpc]{Trivandrum manuscript after correction}
% 	\item[Td]{A transcript of the Trivandrum manuscript}
% 	\item[\Tdac]{A transcript of the Trivandrum manuscript before correction}
% 	\item[\Tdpc]{A transcript of the Trivandrum manuscript after correction}
% 	\item[T]{A reading common in \Tm\ and \Td}
% 	\item[L]{Lahore manuscript}
% 	\item[\Lac]{Lahore manuscript before correction}
% 	\item[\Lpc]{Lahore manuscript after correction}
% 	\item[M]{Madras Government Oriental Manuscript Library manuscript}
% 	\item[\Mac]{Madras Government Oriental Manuscript Library manuscript before correction}
% 	\item[\Mpc]{Madras Government Oriental Manuscript Library manuscript after correction}
% 	\item[A]{Adyar Library manuscript}
% 	\item[\Aac]{Adyar Library manuscript before correction}
% 	\item[\Apc]{Adyar Library manuscript after correction}
% 	\item[\Me]{1952 edition of the YVi}
% 	% \item[L]{}
% 	% \item[L]{}
% 	% \item[L]{}
% 	% \item[L]{}
% \end{description}